\documentclass[10pt]{article}
\usepackage[margin=0.6in]{geometry}
\usepackage{amsmath,amssymb}
\usepackage{multicol}
\usepackage{booktabs}
\usepackage{array}
\usepackage{enumitem}
\usepackage[dvipsnames]{xcolor}

\pagestyle{empty}

% Tableau signs
\newcommand{\T}{\mathbb{T}}
\newcommand{\F}{\mathbb{F}}

% Vertical dots shorthand
\newcommand{\vmark}{\,\vert\,}

% Section formatting
\newcommand{\ruleheader}[1]{\medskip\noindent\textbf{\large #1}\par\smallskip\hrule\smallskip}
\newcommand{\rulesubheader}[1]{\smallskip\noindent\textbf{#1}\par\smallskip}

% Branching notation
\newcommand{\branches}[2]{%
  \begin{tabular}{@{}c@{\qquad}c@{}}
  #1 & #2
  \end{tabular}%
}

\setlength{\parindent}{0pt}
\setlength{\parskip}{2pt}

\begin{document}

\begin{center}
{\LARGE\bfseries Midterm}\\[4pt]
{\large Phil 305 --- Propositional and Modal Logic}
\end{center}

\ruleheader{Part 1: S5}

Use the simplified rules for \textbf{S5} to build tableaux for each of these claims. Say which of them are valid, i.e., are theorems of \textbf{S5}.

\begin{enumerate}
\item $\Box(\Box p \rightarrow q) \vee \Box(\Box q \rightarrow p)$
\item $\Diamond \Box p \rightarrow \Box \Diamond p$
\item $(\neg p \vee \neg \Box \Box p) \rightarrow \neg \Box p$
\end{enumerate}

\ruleheader{Part 2: K}

Use the rules for \textbf{K} to build tableaux for each of these claims. Say which of them are valid, i.e., are theorems of \textbf{K}.

\begin{enumerate}
\setcounter{enumi}{3}
\item $(\Diamond p \wedge \Diamond q) \rightarrow \Box(p \vee q)$
\item $\Diamond (p \rightarrow p) \vee \Diamond \neg (q \rightarrow q)$
\item $(p \wedge \neg \neg \Box q) \rightarrow (\neg \neg p \wedge \Box q)$
\end{enumerate}

\ruleheader{Part 3: S4}

Use the rules for \textbf{S4} to build tableaux for each of these claims. Say which of them are valid, i.e., are theorems of \textbf{S4}. Remember \textbf{S4} is the logic \textbf{KT4}, so it has the rules for \textbf{T} and \textbf{4} as well as \textbf{K}

\begin{enumerate}
\setcounter{enumi}{6}
\item $\Box(p \rightarrow q) \rightarrow (\Diamond \Diamond p \rightarrow \Diamond q)$
\item $(p \wedge \Diamond \Box p) \rightarrow \Box p$
\end{enumerate}

\newpage
\rulesubheader{1. $\Box(\Box p \rightarrow q) \vee \Box(\Box q \rightarrow p)$ in \textbf{S5}}

\newpage
\rulesubheader{2. $\Diamond \Box p \rightarrow \Box \Diamond p$ in \textbf{S5}}

\newpage
\rulesubheader{3. $(\neg p \vee \neg \Box \Box p) \rightarrow \neg \Box p$ in \textbf{S5}}

\newpage
\rulesubheader{4. $(\Diamond p \wedge \Diamond q) \rightarrow \Box(p \vee q)$ in \textbf{K}}

\newpage
\rulesubheader{5. $\Diamond (p \rightarrow p) \vee \Diamond \neg (q \rightarrow q)$ in \textbf{K}}

\newpage
\rulesubheader{6. $(p \wedge \neg \neg \Box q) \rightarrow (\neg \neg p \wedge \Box q)$ in \textbf{K}}

\newpage
\rulesubheader{7. $\Box(p \rightarrow q) \rightarrow (\Diamond \Diamond p \rightarrow \Diamond q)$ in \textbf{S4}}

\newpage
\rulesubheader{8. $(p \wedge \Diamond \Box p) \rightarrow \Box p$ in \textbf{S4}}

\end{document}